\section{Introduction}
\label{sec:introduction}

Long-horizon reasoning is not a longer version of one-shot question answering. At the
start of a mathematical investigation, neither the useful intermediate claims nor
the reachable theorem frontier may be known. An engineering user may express a
solution before understanding the actual performance, reliability, or integration
constraint. In system verification, a failed proof obligation may reveal a defect
in the implementation, an inconsistency in the specification, or a mistaken
assumption shared by both. In each case, execution changes what the user knows
about the problem itself.

This creates a different systems objective from autonomous completion of a fixed
prompt. The system must preserve a standing intent while allowing the operational
goal, constraints, and acceptance criteria to become more precise. It must value
verified intermediate progress---a counterexample, a stronger bound, a measured
bottleneck, a minimal failing trace, or a corrected specification---even when no
terminal artifact yet exists. It must also expose these updates to the user,
whose judgment can revise priorities and constraints for the next mission. The
system challenge is therefore not merely to persist, but to determine when evidence
supports continuing the current route and when it warrants a pivot in approach,
objective, constraints, or verification criteria.

Large language model agents have progressed from tool use and software execution to
complete research loops. ReAct, SWE-agent, and OpenHands established practical
interfaces between models and external environments
~\citep{yao2023react,yang2024sweagent,wang2024openhands}. The AI Scientist,
CycleResearcher, AutoSci, FARS, and Arbor extend this interaction into iterative
research programs covering different combinations of planning, experiments, review,
memory, and scientific
communication~\citep{lu2024aiscientist,weng2025cycleresearcher,qian2026autosci,
tang2026fars,jin2026arbor}. Their progress makes the central systems problem clear:
long execution becomes reasoning capability only when verified results accumulate
across iterations and the runtime can change direction when those results invalidate
the current route or task formulation.

In practice, long-running agents fail in four recurrent ways, consistent with
results showing that agent success declines as software-task horizons grow
~\citep{kwa2025longtasks}. First,
\emph{continuity} is fragile: a growing transcript is eventually compacted or
dropped, and later calls spend time reconstructing decisions that were already
made. Second, \emph{acceptance} is ambiguous when the component that performs the
work also declares it complete. Third, \emph{experience} is lost when only the
final artifact survives; failed branches, contaminated measurements, and useful
procedures disappear with the session that produced them. Hierarchical and
workflow-oriented memory systems address parts of this problem
~\citep{packer2023memgpt,wang2024workflowmemory,xu2025amem}. A larger context window
extends session continuity; durable institutional memory additionally requires
explicit state, ownership, and update rules.
Fourth, the task contract itself can be wrong or incomplete. A runtime that treats
the initial prompt, specification, or evaluator as immutable may optimize the wrong
quantity with increasing efficiency. Conversely, a runtime that freely rewrites
its objective turns learning into unauditable goal drift. The missing mechanism is
verified contract refinement: evidence must identify why a goal or constraint is
insufficient, and an authorized decision must record what changed and what remains
invariant.

We present \systemname (Autonomous Research Generation and Understanding System), a
verification-centered runtime for persistent and pivot-capable long-horizon reasoning. Its
central abstraction is a sequence of \emph{bounded missions} executed against a
durable, revisable task contract.
Four model-driven roles divide authority: a Manager anchors the objective and
campaign state, a Planner selects the next units of work, an Engineer implements
and evaluates them, and a Reviewer inspects artifacts and run outputs before issuing
the completion verdict or identifying a contract-level contradiction. The user
remains outside the repetitive round loop but retains authority at consequential
boundaries. Manager-visible evidence, alternatives, and unresolved questions make
clarification possible when the correct goal or constraint cannot be settled by
artifacts alone. These roles operate across three planes---control,
execution, and records---so that scheduling, work, and record keeping remain
separable even though they participate in one continuous loop.

This report describes the system and evaluations available through August 4, 2026;
historical benchmark traces are reported under the configurations that produced
them.

The runtime is designed to accumulate two coupled forms of state without updating
the base model. First, it self-evolves the \emph{approach} between the current goal
and constraints: each Stage retains attempted routes, measurements, Reviewer
verdicts, accepted and rejected results, reusable skills, verifiers, and routing
decisions, so later missions begin from a changed search policy rather than merely
a longer transcript. Second, verified evidence can sharpen the user-facing task
contract: goals can be decomposed or reframed, constraints can be discovered or
corrected, and acceptance tests can be strengthened. Verifier-gated runtime
self-evolution is the engine that improves execution under the current contract.
Verification-guided pivoting changes that contract only when evidence exposes a
material mismatch; user clarification or authorization governs changes that cannot
be settled by task-native evidence alone.
The term \emph{dense intelligence} describes the corresponding design goal: keep
decision, execution, and verification connected over the research horizon rather
than concentrating useful reasoning into isolated bursts. Section~\ref{sec:problem}
formalizes this conversion and the experiments measure its operating factors.

We sharpen this design language into a verified-progress theory. The final artifact
is a projection of a richer typed trajectory; intermediate states can carry
task-native value; reviewed compression retains the parts needed for later
decisions; and an accepted update counts as compounding only when it reduces future
task risk or resource cost. Contract refinement is admissible only when it preserves
the standing intent, cites new evidence, records the changed assumptions, and passes
the appropriate user or Manager authority boundary. This separates a justified
pivot from goal drift.

We evaluate the resulting system across seven benchmark arenas covering software
repair, GPU kernels, language-model training, training speed, research-assistant
tasks, and mathematical data synthesis. Beyond the headline SWE-Bench result, the
suite includes 76.8\% on AARRI-Bench and a 28.0 mathematical-data gap, both above
their reported references. SWE-Bench Pro is one row in this benchmark suite; its
731-task trajectory additionally enables analyses of Skill/Wiki evolution and
Reviewer intervention. A de-duplicated inventory records 41 research artifacts
across six programs as a separate breadth measure. We additionally report one
mathematical campaign as a representative vertical trace, showing how the four
roles hand off the same objective through recovery, route testing, proof
production, revision, acceptance, and retained state; it preserves one falsified
route and six proof-backed theorem-frontier updates. Finally, we reconstruct six
Argus paper-production campaigns spanning 254 missions and 16 Stage rollbacks from
their structured Stage and role logs. These traces are especially relevant to
pivoting because several initial positive hypotheses fail and become narrower
diagnostic or negative-results contributions through evidence-backed rollback.

This report makes seven contributions:

\begin{enumerate}[label=\textbf{C\arabic*.}]
  \item We formulate long-horizon reasoning as verified progress over a revisable task
    contract $K_t=(\iota,o_t,c_t,v_t)$ comprising standing intent, objectives, constraints, and
    verification criteria, rather than completion against an immutable prompt.
  \item We introduce a role-separated runtime for verification-guided persistence and pivoting that
    combines durable intent, bounded execution, independent review, user-visible
    contract refinement, append-only records, and fixed-model runtime evolution.
  \item We report one unified table across seven benchmark arenas, preserving each
    benchmark's backbone, backend, protocol, native metric, and reference.
  \item We analyze the SWE-Bench Pro row in depth: \systemname reaches approximately
    78\% accuracy versus 59\% for Direct Copilot at 1.41$\times$ aggregate Tokens;
    mature Waves are more efficient, and Reviewer intervention yields 34 verifier
    recoveries and 22 strict review-loop rescues.
  \item We reconstruct one mathematical campaign as a role-resolved Agent trace,
    showing why falsified routes and intermediate theorem-frontier updates are
    valuable before a terminal proof and how they reshape subsequent work.
  \item We present a six-project autonomous paper-production case study covering
    \PaperCaseCampaignHours{} campaign-hours, \PaperCaseRounds{} Engineer rounds,
    \PaperCaseContinues{} Reviewer revisions, and \PaperCaseRollbacks{} Stage
    rollbacks, with all six canonical pipelines reaching submission completion. A
    representative campaign uses seven verified no-go decisions to reframe a failed
    method search as a 4,500-row negative-results study, then repairs two late
    submission defects without discarding the accumulated research state.
  \item We document two artifact-level outcomes: a human-reviewed
    TileLang kernel merged into Flash Linear Attention and the operator-reported
    ACE-1 NPU prototype, while keeping their distinct evidence levels explicit.
\end{enumerate}

The remainder of the report follows a conventional research-paper organization.
Section~\ref{sec:related} positions the system against prior agent and autonomous
research work. Section~\ref{sec:problem} formalizes the operating problem;
Section~\ref{sec:method} presents the runtime; Sections~\ref{sec:methodology} and
\ref{sec:results} describe the evaluation protocol and results; and
Sections~\ref{sec:discussion}--\ref{sec:conclusion} discuss interpretation,
limitations, and conclusions. The appendix contains detailed experimental tables
and notation.
